\chap{自主团队组织：架构分工、递归委派与长期经验}
\label{sec:organization}
\subsection{长期职责与临时任务}
开发系统中的一些工作具有持续性，例如维护任务状态语义、执行器接口或用户交互。另一些工作具有临时性，例如完成一次跨模块改进、调查偶发故障或比较两个方案。组织时需要同时容纳这两种时间尺度。持续成员围绕稳定职责积累经验，临时成员则围绕当前任务聚集必要背景和协调关系。

长期沟通入口可以由 Leader 承担。它持续理解人的方向、偏好和项目承诺，并把工作委派给合适成员。Leader 也可以委派复杂任务的管理、验证跟踪和交付组织。若每次跨模块问题都必须回到顶层进行细节决策，组织会出现串行瓶颈，其上下文也会快速被执行信息占据。

因此，责任可以递归委派。任务负责人获得具体目标、约束和相关成员入口，再自行组织实现与验证。返回上层的信息聚焦成果、对原目标的影响和需要整体取舍的问题。细节继续保存并可按需访问。委派完成的依据是可核查的结果和明确的未完事项，而非成员数量或消息流量。

\subsection{沿架构职责组织成员}
Portal 的实践讨论倾向于沿功能、服务或包所体现的职责组织成员，以保持模块经验和接口连续性。与固定的需求、实现、测试阶段角色相比，这种方式可能减少同一背景在多个成员之间反复重建的成本。它是一项基于具体实践提出的组织假设，其适用性需要在不同项目中验证。

模块边界应依据职责与变化影响判断。一个目录包含若干文件，并不自动意味着它是合适的持续责任范围。可以观察某类变化通常由哪些知识决定、哪些接口会受到影响，以及是否经常迫使无关成员理解内部细节。若多个成员总要共同修改同一个内部决定，边界可能需要调整；若一个成员维护过多相互独立的领域，则拆分可能有益。

经典模块化研究将隐藏可能变化的设计决定作为分解依据之一\cite{parnas1972}。对 Agent 团队而言，职责边界还决定上下文应积累在哪里。执行器成员熟悉停止信号、资源释放和异常恢复；状态服务成员熟悉请求语义、幂等和终态判断。清楚的协议允许双方在局部深入，同时把跨边界的共同假设显式化。

\fig{04-organization}{组织关系与软件关系分别描述。持续成员关联维护范围，临时负责人组织跨模块任务；组织图中的委派边与软件运行时依赖具有不同含义。}{organization}

团队图与软件架构图可以关联展示，但不必严格一一对应。一个成员可能维护多个小模块，一个复杂模块也可能需要临时协作。通信频率高也不一定意味着架构有问题：一次协议升级可能合理地需要密集沟通。判断应联系变化类型、重复发生的协调需求及维护成本。

\subsection{协议是自主协作的共同依据}
模块之间需要明确可依赖的行为，包括输入、输出、状态、失败含义和兼容条件。协议应同时支持实现与解释。以取消为例，任务服务向执行器发送停止意图，执行器返回接收信息，随后报告实际终态。若各方对“接收”与“完成”的理解不同，单个模块测试都可能通过，组合起来却误导用户。

协议修改应联系具体调用方和使用任务。负责人可以先构造典型与异常交互，再让相关成员核对自身承担的部分。对于跨版本变化，还需要确定短暂共存期间的解释方式，例如旧执行器只提供一个状态时，界面如何避免过度承诺。接口稳定性并不要求协议永远不变，而是要求变化的条件和影响能够被理解与处理。

跨模块协调的结果应足够简洁，可被后续工作重新使用。记录共同决定、适用版本、重要例外和验证入口，往往比保存全部讨论更有用。若多个临时任务同时修改同一契约，持续负责人可以帮助统一解释，避免出现相互竞争的局部规则。

\subsection{上下文分配与递归委派}
上下文是稀缺的执行资源。它既占用推理成本，也影响成员注意什么信息。对复杂任务，合理的做法是把与局部判断有关的材料交给负责成员，把稳定结论、重要依赖和结果入口返回上层。所有成员接收所有日志，会使组织丧失分工带来的信息优势。

fork 可以作为派生任务上下文的一种手段，使新成员继承分出时的相关背景，再深入局部工作。但继承也可能携带过时假设和无关材料。派生时应明确目标、边界、所依据的状态以及需要返回什么结果。成员在分支中发现新事实后，不能假定父成员会自动知道，需要通过合适的结果关联更新共享理解。

委派的收益可以粗略理解为节省的深入调查与并行等待，减去交接、协调和重新整合的成本。对于很小的修改，继续委派可能更慢；对于需要大量局部探索的任务，独立上下文更有价值。Leader 或任务负责人应能够改变任务粒度，并根据过去的实际成本调整组织。

\fig{05-context}{上下文与经验的生命周期。持久材料按任务进入当前上下文，执行结果经过核对后提炼为带适用条件的经验，供下一次任务和会话恢复使用。}{context}

\subsection{连续身份、恢复与经验更新}
持续成员的身份应与单次模型会话分离。模型更新、上下文压缩或执行中断之后，成员仍应能够找到自己的职责、当前承诺、重要决定和未完工作。恢复所需的信息应短而有结构，并保留深入材料的入口。只有一个“继续上次任务”的模糊摘要，往往不足以恢复实际状态。

经验可以分为对模块的稳定认识、具体问题的处理案例和可执行工具。稳定认识帮助快速建立背景，案例帮助识别相似问题，工具则将确定性操作固化为可重复能力。三者的维护方式不同。一个已失效的接口描述需要更新，一次罕见故障的历史证据仍可保留，一个脚本则需要随接口变化进行验证。

Reflexion 探索通过语言反馈和记忆影响后续尝试\cite{reflexion}，Generative Agents 研究记忆、反思和规划在持续行为中的作用\cite{generativeagents}。这些研究提示外部经验可能改善后续行动。软件开发中的经验还需要与版本、契约和运行条件关联，其长期收益应通过后续维护任务检验。

成员退出时，成果与重要判断应能够被继任者使用。退出不应导致任务只有原成员才知道在哪里，也不应把所有会话历史无差别转交。可以以职责为单位保留当前契约、已知问题、活跃任务和证据入口；继任者在真实任务中逐步补充必要细节。

\subsection{并行工作中的协调边界}
成员自主工作时，某些局部决定可能同时触及共享契约。协调应集中在这个共同边界，而不要求所有成员汇报全部内部过程。例如，界面成员准备合并两个状态，执行器成员则准备增加一种停止异常；双方需要先核对状态的用户含义，其余内部实现可以继续并行。

当两项任务竞争同一资源或改变同一承诺时，可以由临时负责人组织取舍。它需要了解相关目标、依赖和时间要求，而不必重新掌握全部模块实现。取舍完成之后，将共同结论返回责任成员，并保留影响范围。这个过程有助于避免顶层 Leader 成为每次局部冲突的中转站。

组织还应能识别工作停滞。成员可能已经完成局部实现，却在等待另一方确认协议；也可能因目标变化而继续执行已经失去价值的任务。共享状态应表达等待什么、由谁推进和何时需要重新判断。对人展示时可以进一步压缩为实际阻塞与可选行动，内部则保留足够的责任联系。

\subsection{组织调整与独立核对}
按模块组织不妨碍专业评审或独立验证。实现者通常掌握最多局部知识，但也可能沿用自己的错误假设。对跨模块行为、复杂协议或重要用户任务，可以让另一成员从实际入口核对结果。独立性应体现在信息来源与判断方式，而不只是角色名称不同。

组织本身也需要接受评价。重复调查、无效交接、等待顶层决定和共享契约冲突，都是可以观察的信号。系统可以据此合并职责、拆分过大范围或设置临时协调者。组织调整应解释要解决的具体问题，并观察之后是否减少成本，避免通过不断增加角色来处理角色之间产生的问题。

Conway 对组织沟通与系统设计关系的讨论\cite{conway1968}提醒我们，协作结构会影响软件边界。对于可动态创建成员的系统，这种关系提供了主动调整的机会。机会的实现仍依赖对实际变化的理解：组织图的整齐并不能替代软件内部职责的清晰。
